\documentclass[12pt]{article}
\usepackage[utf8]{inputenc}
\usepackage{amsmath,amssymb,latexsym}
\usepackage{enumerate,comment}
\usepackage{graphicx}
\usepackage[spanish]{babel}
\usepackage{xypic}

\textheight=25cm \textwidth=18cm \topmargin=-2cm \oddsidemargin=-1cm

\newcommand{\abs}[1]{\left\vert#1\right\vert}
\newcommand{\norm}[1]{\left\Vert#1\right\Vert}

\newcommand{\Co}{\mathfrak C}
\newcommand{\diag}{\mathrm{diag}}
\newcommand{\Ext}{\mathrm{ext}}
\newcommand{\F}{\mathfrak F}
\newcommand{\h}{\mathrm h}
\newcommand{\Int}{\mathrm{int}}
\newcommand{\n}{\mathfrak{n}}
\newcommand{\N}{\mathbb N}
\newcommand{\Q}{\mathbb Q}
\newcommand{\R}{\mathbb R}
\newcommand{\US}{\mathbb{S}^2}
\newcommand{\x}{\mathrm x}
\newcommand{\xk}{(\x_k)_{k\in \N}}
\newcommand{\y}{\mathrm y}
\newcommand{\z}{\mathrm z}
\newcommand{\Z}{\mathbb Z}

\newcommand{\eps}{\varepsilon}
\newcommand{\lbil}{\left <}
\newcommand{\rbil}{\right >}
\newcommand{\To}{\longrightarrow}
\newcommand{\mTo}{\longmapsto}

\begin{document}
\pagestyle{empty}
\begin{center}
{\textsc{Cálculo Diferencial e Integral III} \\
\medskip 
\large{Ejercicios I - 1 \\
(El espacio euclidiano de dimensión $n$)}}
\end{center}
\bigskip

En general, $\Omega,\Delta, \dots$ serán subconjuntos del espacio euclidiano $\R^n$.

\begin{enumerate}

\item 
\begin{enumerate}[(1)]
\item Verifica que las funciones
\[ 
\norm{\x}_{\max} := \max\{\abs{x_1},\ldots,\abs{x_n}\} \quad \text{y}\quad 
\norm{\x}_1 := \abs{x_1} + \cdots + \abs{x_n} 
\]
satisfacen las propiedades de una norma en $\R^n$.

\textbf{Solución: } Veamos que \(\norm{\x}_{\max}\) es una norma:
    \begin{itemize}
        \item \textbf{Positividad:} Por definición del valor absoluto, \(\abs{x_i} \geq 0\) para toda \(i \in [n]\), por lo que \(\max\{\abs{x_1}, \dots, \abs{x_n}\} \geq 0\). Así, \(\norm{\x}_{\max} \geq 0\) para toda \(\x \in \R^n\).
        
        Además, \(\norm{\x}_{\max} = 0\) si y solo si \(\max\{\abs{x_1}, \dots, \abs{x_n}\} = 0\), lo que ocurre si y solo si \(\abs{x_i} = 0\) para toda \(i \in [n]\), es decir, si y solo si \(\x = 0\).
        
        \item \textbf{Homogeneidad:} Sea \(\lambda \in \R\). Usando la propiedad del valor absoluto \(\abs{\lambda x_i} = \abs{\lambda} \abs{x_i}\), tenemos:
        \[
        \norm{\lambda \x}_{\max} = \max_{1 \leq i \leq n}\{\abs{\lambda x_i}\} = \max_{1 \leq i \leq n}\{\abs{\lambda} \abs{x_i}\} = \abs{\lambda} \max_{1 \leq i \leq n}\{\abs{x_i}\} = \abs{\lambda} \norm{\x}_{\max}.
        \]
        
        \item \textbf{Desigualdad del triángulo:} Por definición, 
        \[
        \norm{\x}_{\max} + \norm{\y}_{\max} = \max_{i\in[n]}\{\abs{x_i}\} + \max_{i\in[n]}\{\abs{y_i}\}.
        \]
        Para cualquier \(j \leq n\), se tiene que \(\max_{i\in[n]}\{\abs{x_i}\} \geq \abs{x_j}\) y \(\max_{i\in[n]}\{\abs{y_i}\} \geq \abs{y_j}\). Entonces,
        \[
        \max_{i\in[n]}\{\abs{x_i}\} + \max_{i\in[n]}\{\abs{y_i}\} \geq \abs{x_j} + \abs{y_j} \geq \abs{x_j + y_j}
        \]
        para cualquier \(j \leq n\). Esto implica que
        \[
        \max_{i\in[n]}\{\abs{x_i}\} + \max_{i\in[n]}\{\abs{y_i}\} \geq \max_{i \in [n]}\{\abs{x_j + y_j}\},
        \]
        y por lo tanto, \(\norm{\x}_{\max} + \norm{\y}_{\max} \geq \norm{\x + \y}_{\max}\).
    \end{itemize}
    
Ahora, veamos que \(\norm{\x}_1\) es una norma:
\begin{itemize}
    \item \textbf{Positividad:} Por definición del valor absoluto, \(\abs{x_i} \geq 0\) y ya que la suma de valores no negativos es no negativa, \(\norm{\x}_1 \geq 0\) para cualquier vector \(\x \in \R^n\).
    
    Además, ya que \(\sum_{i \in [n]}\abs{x_i} = 0\) si y solo si \(\abs{x_i} = 0\) para cualquier índice \(i\), se sigue que \(\norm{\x}_1 = 0\) si y solo si \(\x = 0\).
    
    \item \textbf{Homogeneidad:} Sea \(\lambda \in \mathbb{R}\). Por definición del producto escalar, tenemos que 
    \[
    \norm{\lambda \x}_1 = \sum_{i = 1}^{n}\abs{\lambda x_i} = \sum_{i = 1}^{n}\abs{\lambda}\abs{x_i} = \abs{\lambda}\sum_{i = 1}^{n}\abs{x_i} = \abs{\lambda}\norm{\x}_1.
    \]
    
    \item \textbf{Desigualdad del triángulo:} Reescribimos
    \[
    \norm{\x}_1 + \norm{\y}_1 = \sum_{i = 1}^{n}\abs{x_i} + \sum_{i = 1}^{n}\abs{y_i}.
    \]
    Sumando término a término \(\abs{x_i} + \abs{y_i}\), obtenemos la expresión
    \[
    \sum_{i = 1}^{n}\abs{x_i} + \sum_{i = 1}^{n}\abs{y_i} = \sum_{i = 1}^{n}(\abs{x_i} + \abs{y_i}) \geq \sum_{i = 1}^{n}\abs{x_i + y_i}.
    \]
    Por lo tanto, 
    \[
    \norm{\x}_1 + \norm{\y}_1 \geq \norm{\x + \y}_1.
    \]
\end{itemize}

\item Describe las bolas unitarias en $\R^2$ con respecto a las normas $\norm{\x}_{\max}$ y $\norm{\x}_1$. 
\begin{itemize}
    \item La bola unitaria es el conjunto \(B_{\max} = \{(x_1, x_2) \in \R^2 : \norm{\x}_{\max} \leq 1\}\), lo que significa que
    \[
    \max\{\abs{x_1}, \abs{x_2}\} \leq 1
    \]
    que es equivalente a
    \[
    -1 \leq x_1 \leq 1 \quad \text{y} \quad -1 \leq x_2 \leq 1,
    \]
    por lo que la bola unitaria es un cuadrado de lado 2 centrado alrededor del origen:
    \[
    B_{\max} = [-1, 1] \times [-1, 1].
    \]
    \item Para la norma \(\norm{\x}_1\), la bola unitaria se define como el conjunto \(B_1 = \{(x_1, x_2) \in \R^2 : \norm{\x}_1 \leq 1\}\). Es decir, 
    \[
    \abs{x_1} + \abs{x_2} \leq 1.
    \]
    Esto describe un rombo centrado en el origen con vértices en los puntos \((1, 0)\), \((0, 1)\), \((-1, 0)\) y \((0, -1)\).
\end{itemize}
\end{enumerate}

\item Demuestra que se cumplen las siguientes desigualdades:
\begin{enumerate}[(1)]
\item $\norm{\x}_{\max}\leq \norm{\x}\leq \sqrt{n}\,\norm{\x}_{\max}$.
\item $\frac{1}{\sqrt{n}}\norm{\x}_1\leq \norm{\x}\leq \norm{\x}_1$.
\end{enumerate}

\textbf{Solución: } Recordemos que la norma euclidiana se define como
\[
\norm{\x} = \sqrt{x_1^2 + x_2^2 + \cdots + x_n^2}.
\]
Demostraremos cada desigualdad por separado.

\begin{enumerate}[(1)]
\item \textbf{Primera desigualdad:} $\norm{\x}_{\max}\leq \norm{\x}$.

Sea $k \in \{1, \dots, n\}$ un índice tal que $\abs{x_k} = \max\{\abs{x_1}, \dots, \abs{x_n}\} = \norm{\x}_{\max}$. Entonces,
\[
\norm{\x}_{\max}^2 = \abs{x_k}^2 = x_k^2 \leq \sum_{i=1}^{n} x_i^2 = \norm{\x}^2.
\]
Tomando raíz cuadrada en ambos lados (ambas cantidades son no negativas), obtenemos
\[
\norm{\x}_{\max} \leq \norm{\x}.
\]

\textbf{Segunda desigualdad:} $\norm{\x}\leq \sqrt{n}\,\norm{\x}_{\max}$.

Para cada $i \in \{1, \dots, n\}$, se tiene que $\abs{x_i} \leq \norm{\x}_{\max}$, por lo que $x_i^2 = \abs{x_i}^2 \leq \norm{\x}_{\max}^2$. Sumando sobre todos los índices $i$, obtenemos
\[
\norm{\x}^2 = \sum_{i=1}^{n} x_i^2 \leq \sum_{i=1}^{n} \norm{\x}_{\max}^2 = n\,\norm{\x}_{\max}^2.
\]
Nuevamente, tomando raíz cuadrada en ambos lados, concluimos que
\[
\norm{\x} \leq \sqrt{n}\,\norm{\x}_{\max}.
\]
Por lo tanto, combinando ambas desigualdades, se cumple
\[
\norm{\x}_{\max}\leq \norm{\x}\leq \sqrt{n}\,\norm{\x}_{\max}.
\]

\item \textbf{Primera desigualdad:} $\frac{1}{\sqrt{n}}\norm{\x}_1\leq \norm{\x}$.

Esta desigualdad es una consecuencia directa de la desigualdad de Cauchy-Schwarz. Consideremos los vectores $\x = (x_1, \dots, x_n)$ y $\mathbf{1} = (1, 1, \dots, 1)$ en $\R^n$. Aplicando la desigualdad de Cauchy-Schwarz, tenemos:
\[
\abs{\lbil \x, \mathbf{1} \rbil} \leq \norm{\x}\,\norm{\mathbf{1}}.
\]
Calculamos cada término:
\[
\abs{\lbil \x, \mathbf{1} \rbil} = \abs{\sum_{i=1}^{n} x_i \cdot 1} = \abs{\sum_{i=1}^{n} x_i} \leq \sum_{i=1}^{n} \abs{x_i} = \norm{\x}_1.
\]
y
\[
\norm{\mathbf{1}} = \sqrt{1^2 + 1^2 + \cdots + 1^2} = \sqrt{n}.
\]
Por lo tanto, la desigualdad de Cauchy-Schwarz nos da
\[
\norm{\x}_1 = \sum_{i=1}^{n} \abs{x_i} \leq \norm{\x}\,\sqrt{n}.
\]
Dividiendo ambos lados entre $\sqrt{n}$ (que es positivo), obtenemos
\[
\frac{1}{\sqrt{n}}\norm{\x}_1 \leq \norm{\x}.
\]

\textbf{Segunda desigualdad:} $\norm{\x}\leq \norm{\x}_1$.

Demostraremos esta desigualdad elevando al cuadrado ambos lados. Como ambas cantidades son no negativas, es suficiente demostrar que $\norm{\x}^2 \leq \norm{\x}_1^2$. Tenemos:
\[
\norm{\x}^2 = \sum_{i=1}^{n} x_i^2.
\]
Por otro lado,
\[
\norm{\x}_1^2 = \left(\sum_{i=1}^{n} \abs{x_i}\right)^2 = \sum_{i=1}^{n} \abs{x_i}^2 + 2\sum_{1 \leq i < j \leq n} \abs{x_i}\abs{x_j}.
\]
Observamos que $\sum_{i=1}^{n} \abs{x_i}^2 = \sum_{i=1}^{n} x_i^2 = \norm{\x}^2$ y que el término $2\sum_{1 \leq i < j \leq n} \abs{x_i}\abs{x_j}$ es no negativo (es una suma de productos de valores absolutos, todos $\geq 0$). Por lo tanto,
\[
\norm{\x}_1^2 = \norm{\x}^2 + 2\sum_{1 \leq i < j \leq n} \abs{x_i}\abs{x_j} \geq \norm{\x}^2.
\]
Tomando raíz cuadrada en ambos lados, concluimos que
\[
\norm{\x} \leq \norm{\x}_1.
\]
Por lo tanto, combinando ambas desigualdades, se cumple
\[
\frac{1}{\sqrt{n}}\norm{\x}_1\leq \norm{\x}\leq \norm{\x}_1.
\]
\end{enumerate}

\item Describe el interior de
\[ \Omega = \{ (x,y,z)\in \R^3 : 0\leq x < 1, y^2+z^2\leq 1\}. \]

\textbf{Solución: } Recordemos que un punto $\mathbf{p} \in \Omega$ pertenece al interior de $\Omega$, denotado por $\Int(\Omega)$, si existe una bola abierta centrada en $\mathbf{p}$ completamente contenida en $\Omega$. Es decir, si existe $\varepsilon > 0$ tal que $B(\mathbf{p}, \varepsilon) \subset \Omega$.

Analicemos las condiciones que definen a $\Omega$:
\begin{itemize}
    \item La condición $0 \leq x < 1$ restringe la primera coordenada.
    \item La condición $y^2 + z^2 \leq 1$ restringe las otras dos coordenadas a un disco cerrado de radio 1 en el plano $yz$.
\end{itemize}

El conjunto $\Omega$ es como un ``cilindro sólido'' en $\R^3$, cuyo eje es el eje $x$, con radio 1, pero con tapas en $x=0$ y $x=1$. La tapa en $x=0$ está incluida (debido al signo $\leq$), mientras que la tapa en $x=1$ no está incluida (debido al signo $<$). La superficie lateral del cilindro, definida por $y^2+z^2 = 1$, sí está incluida.

Para que un punto $(x,y,z) \in \Omega$ esté en el interior, debe ser posible moverse un poco en \textbf{cualquier} dirección sin salirse de $\Omega$. Esto nos lleva a las siguientes condiciones:

\begin{enumerate}
    \item \textbf{Condición sobre $x$:} 
    \begin{itemize}
        \item Si $x = 0$, cualquier bola centrada en el punto contendrá puntos con primera coordenada negativa (es decir, $x < 0$), los cuales no están en $\Omega$. Por lo tanto, los puntos con $x=0$ no pueden estar en el interior.
        \item Si $0 < x < 1$, podemos elegir un $\varepsilon > 0$ lo suficientemente pequeño (por ejemplo, $\varepsilon < \min\{x, 1-x\}$) de modo que si $|x' - x| < \varepsilon$, entonces $0 < x' < 1$. Así, la condición sobre $x$ no es un obstáculo para estar en el interior.
    \end{itemize}
    
    \item \textbf{Condición sobre $y,z$:}
    \begin{itemize}
        \item Si $y^2 + z^2 = 1$ (el punto está en la superficie lateral del cilindro), cualquier bola centrada en el punto contendrá puntos con $y'^2 + z'^2 > 1$, los cuales no están en $\Omega$. Por lo tanto, estos puntos no pueden estar en el interior.
        \item Si $y^2 + z^2 < 1$, podemos elegir un $\varepsilon > 0$ lo suficientemente pequeño de modo que si $(y',z')$ está cerca de $(y,z)$, entonces $y'^2 + z'^2 < 1$. Así, la condición sobre $y,z$ no es un obstáculo para estar en el interior.
    \end{itemize}
\end{enumerate}

Combinando ambas condiciones, concluimos que los puntos del interior deben cumplir estrictamente ambas desigualdades. Por lo tanto,
\[
\Int(\Omega) = \{ (x,y,z)\in \R^3 : 0 < x < 1, \ y^2+z^2 < 1 \}.
\]

Este conjunto es un cilindro abierto de radio 1, con eje en el eje $x$, y con tapas abiertas en $x=0$ y $x=1$.

\item Encuentra la frontera de 
\[ \Omega = \{ (x,y)\in \R^2 : x^2 - y^2 > 1 \}. \]

\textbf{Solución: }
Consideremos la función $f(x,y) = x^2 - y^2$. Esta función es continua en todo $\R^2$. El conjunto $\Omega$ está definido por la condición $f(x,y) > 1$.

Primero, determinamos el interior y la cerradura. Dado que $f$ es continua, el conjunto $\Omega = f^{-1}((1, \infty))$ es abierto. Por lo tanto, $\Int(\Omega) = \Omega$. Es decir,
\[
\Int(\Omega) = \{ (x,y)\in \R^2 : x^2 - y^2 > 1 \}.
\]

Para encontrar la cerradura, observamos que si tomamos puntos $(x,y)$ que cumplen $x^2 - y^2 > 1$ y los acercamos a la curva $x^2 - y^2 = 1$, podemos aproximarnos tanto como queramos. Por lo tanto, la cerradura de $\Omega$ es el conjunto donde $x^2 - y^2 \geq 1$. Es decir,
\[
\overline{\Omega} = \{ (x,y)\in \R^2 : x^2 - y^2 \geq 1 \}.
\]
Esto se debe a que la adherencia de un conjunto abierto definido por una desigualdad estricta con una función continua es el conjunto definido por la desigualdad no estricta correspondiente.

Ahora, determinamos la frontera. Usando la definición de frontera como $\partial \Omega = \overline{\Omega} \setminus \Int(\Omega)$, tenemos:
\[
\partial \Omega = \{ (x,y)\in \R^2 : x^2 - y^2 \geq 1 \} \setminus \{ (x,y)\in \R^2 : x^2 - y^2 > 1 \}.
\]
La diferencia entre estos dos conjuntos es precisamente el conjunto de puntos donde se cumple la igualdad. Por lo tanto,
\[
\partial \Omega = \{ (x,y)\in \R^2 : x^2 - y^2 = 1 \}.
\]

\item Decir si las siguientes propiedades son ciertas. Prueba tu respuesta.
\begin{enumerate}[(1)]
\item $\Int(\Omega\cup \Delta) = \Int(\Omega)\cup \Int(\Delta),\; \Int(\Omega\cap \Delta) = \Int(\Omega)\cap \Int(\Delta)$.
\item $\Ext(\Omega\cup \Delta) = \Ext(\Omega)\cup \Ext(\Delta),\; \Ext(\Omega\cap \Delta) = \Ext(\Omega)\cap \Ext(\Delta)$.
\end{enumerate}

\textbf{Solución: } Analicemos cada propiedad por separado.

\begin{enumerate}[(1)]
\item \textbf{Propiedades del interior.}

\textbf{Primera propiedad:} $\Int(\Omega\cup \Delta) = \Int(\Omega)\cup \Int(\Delta)$.

Esta propiedad es \textbf{falsa} en general. Solo se cumple una contención.

\textit{Contención verdadera:} $\Int(\Omega)\cup \Int(\Delta) \subseteq \Int(\Omega\cup \Delta)$.

\textit{Demostración:} Si $\x \in \Int(\Omega)$, entonces existe $\varepsilon > 0$ tal que $B(\x, \varepsilon) \subseteq \Omega \subseteq \Omega \cup \Delta$. Por lo tanto, $\x \in \Int(\Omega \cup \Delta)$. Análogamente, si $\x \in \Int(\Delta)$, entonces $\x \in \Int(\Omega \cup \Delta)$. Así, $\Int(\Omega)\cup \Int(\Delta) \subseteq \Int(\Omega\cup \Delta)$.

\textit{Contraejemplo para la otra contención:} Tomemos $\Omega = [0,1]$ y $\Delta = [1,2]$ en $\R$. Entonces $\Omega \cup \Delta = [0,2]$, y tenemos:
\[
\Int(\Omega \cup \Delta) = \Int([0,2]) = (0,2),
\]
mientras que
\[
\Int(\Omega) \cup \Int(\Delta) = \Int([0,1]) \cup \Int([1,2]) = (0,1) \cup (1,2).
\]
Observamos que $1 \in (0,2) = \Int(\Omega \cup \Delta)$, pero $1 \notin (0,1) \cup (1,2) = \Int(\Omega) \cup \Int(\Delta)$. Por lo tanto, $\Int(\Omega \cup \Delta) \not\subseteq \Int(\Omega) \cup \Int(\Delta)$, y la igualdad es falsa.

\textbf{Segunda propiedad:} $\Int(\Omega\cap \Delta) = \Int(\Omega)\cap \Int(\Delta)$.

Esta propiedad es \textbf{verdadera}.

\textit{Demostración:} Probaremos la doble contención.

($\subseteq$) Sea $\x \in \Int(\Omega \cap \Delta)$. Entonces existe $\varepsilon > 0$ tal que $B(\x, \varepsilon) \subseteq \Omega \cap \Delta$. Esto significa que $B(\x, \varepsilon) \subseteq \Omega$ y $B(\x, \varepsilon) \subseteq \Delta$. Por lo tanto, $\x \in \Int(\Omega)$ y $\x \in \Int(\Delta)$, lo que implica que $\x \in \Int(\Omega) \cap \Int(\Delta)$.

($\supseteq$) Sea $\x \in \Int(\Omega) \cap \Int(\Delta)$. Entonces $\x \in \Int(\Omega)$ y $\x \in \Int(\Delta)$. Por lo tanto, existen $\varepsilon_1 > 0$ y $\varepsilon_2 > 0$ tales que $B(\x, \varepsilon_1) \subseteq \Omega$ y $B(\x, \varepsilon_2) \subseteq \Delta$. Sea $\varepsilon = \min\{\varepsilon_1, \varepsilon_2\} > 0$. Entonces $B(\x, \varepsilon) \subseteq B(\x, \varepsilon_1) \subseteq \Omega$ y $B(\x, \varepsilon) \subseteq B(\x, \varepsilon_2) \subseteq \Delta$. Por lo tanto, $B(\x, \varepsilon) \subseteq \Omega \cap \Delta$, lo que implica que $\x \in \Int(\Omega \cap \Delta)$.

Así, $\Int(\Omega\cap \Delta) = \Int(\Omega)\cap \Int(\Delta)$.

\item \textbf{Propiedades del exterior.}

Recordemos que el exterior de un conjunto $\Omega$ se define como $\Ext(\Omega) = \Int(\R^n \setminus \Omega)$. Es decir, el exterior de $\Omega$ es el interior de su complemento.

\textbf{Primera propiedad:} $\Ext(\Omega\cup \Delta) = \Ext(\Omega)\cup \Ext(\Delta)$.

Esta propiedad es \textbf{falsa} en general.

\textit{Contraejemplo:} Tomemos $\Omega = (-\infty, 0]$ y $\Delta = [0, \infty)$ en $\R$. Entonces $\Omega \cup \Delta = \R$. El exterior de $\R$ es vacío:
\[
\Ext(\Omega \cup \Delta) = \Ext(\R) = \Int(\R^n \setminus \R) = \Int(\emptyset) = \emptyset.
\]
Por otro lado,
\[
\Ext(\Omega) = \Ext((-\infty, 0]) = \Int((0, \infty)) = (0, \infty),
\]
\[
\Ext(\Delta) = \Ext([0, \infty)) = \Int((-\infty, 0)) = (-\infty, 0).
\]
Entonces,
\[
\Ext(\Omega) \cup \Ext(\Delta) = (-\infty, 0) \cup (0, \infty) = \R \setminus \{0\} \neq \emptyset.
\]
Por lo tanto, $\Ext(\Omega\cup \Delta) \neq \Ext(\Omega)\cup \Ext(\Delta)$, y la igualdad es falsa.

\textbf{Segunda propiedad:} $\Ext(\Omega\cap \Delta) = \Ext(\Omega)\cap \Ext(\Delta)$.

Esta propiedad es \textbf{falsa} en general.

\textit{Contraejemplo:} Tomemos $\Omega = (-\infty, 1)$ y $\Delta = (0, \infty)$ en $\R$. Entonces $\Omega \cap \Delta = (0, 1)$. Calculamos los exteriores:
\[
\Ext(\Omega \cap \Delta) = \Ext((0,1)) = \Int(\R \setminus (0,1)) = \Int((-\infty, 0] \cup [1, \infty)) = (-\infty, 0) \cup (1, \infty).
\]
Por otro lado,
\[
\Ext(\Omega) = \Ext((-\infty, 1)) = \Int([1, \infty)) = (1, \infty),
\]
\[
\Ext(\Delta) = \Ext((0, \infty)) = \Int((-\infty, 0]) = (-\infty, 0).
\]
Entonces,
\[
\Ext(\Omega) \cap \Ext(\Delta) = (1, \infty) \cap (-\infty, 0) = \emptyset.
\]
Claramente, $\Ext(\Omega \cap \Delta) = (-\infty, 0) \cup (1, \infty) \neq \emptyset = \Ext(\Omega) \cap \Ext(\Delta)$. Por lo tanto, la igualdad es falsa.

\textbf{Observación:} Usando las leyes de De Morgan, se puede demostrar que las propiedades correctas para el exterior son:
\[
\Ext(\Omega \cup \Delta) = \Ext(\Omega) \cap \Ext(\Delta),
\]
\[
\Ext(\Omega \cap \Delta) = \Ext(\Omega) \cup \Ext(\Delta).
\]
Es decir, el exterior intercambia uniones por intersecciones y viceversa, de manera análoga al complemento.
\end{enumerate}

\item Decir si las siguientes propiedades son ciertas. Prueba tu respuesta.
\begin{enumerate}[(1)]
\item $\partial(\Omega\cup \Delta) = \partial \Omega\cup \partial \Delta,\; 
\partial(\Omega\cap \Delta) = \partial \Omega\cap \partial \Delta$.
\item Si $\Omega\subset \Delta$ entonces $\Int(\Omega)\subset \Int(\Delta),\; \Ext(\Omega)\subset \Ext(\Delta),\; \partial \Omega\subset \partial \Delta$.
\end{enumerate}

\textbf{Solución: } Analicemos cada propiedad por separado.

\begin{enumerate}[(1)]
\item \textbf{Propiedades de la frontera.}

\textbf{Primera propiedad:} $\partial(\Omega\cup \Delta) = \partial \Omega\cup \partial \Delta$.

Esta propiedad es \textbf{falsa} en general. Solo se cumple una contención.

\textit{Contención verdadera:} $\partial \Omega \cup \partial \Delta \subseteq \partial(\Omega \cup \Delta)$.

\textit{Demostración:} Sea $\x \in \partial \Omega$. Entonces toda bola abierta centrada en $\x$ contiene puntos de $\Omega$ y puntos de $\R^n \setminus \Omega$. Como $\Omega \subseteq \Omega \cup \Delta$, toda bola centrada en $\x$ contiene puntos de $\Omega \cup \Delta$. Además, como $\R^n \setminus (\Omega \cup \Delta) \subseteq \R^n \setminus \Omega$, toda bola centrada en $\x$ contiene puntos de $\R^n \setminus (\Omega \cup \Delta)$. Por lo tanto, $\x \in \partial(\Omega \cup \Delta)$. Análogamente, si $\x \in \partial \Delta$, entonces $\x \in \partial(\Omega \cup \Delta)$. Así, $\partial \Omega \cup \partial \Delta \subseteq \partial(\Omega \cup \Delta)$.

\textit{Contraejemplo para la otra contención:} Tomemos $\Omega = [0,1]$ y $\Delta = [1,2]$ en $\R$. Entonces $\Omega \cup \Delta = [0,2]$. Calculamos las fronteras:
\[
\partial(\Omega \cup \Delta) = \partial([0,2]) = \{0, 2\},
\]
mientras que
\[
\partial \Omega \cup \partial \Delta = \partial([0,1]) \cup \partial([1,2]) = \{0, 1\} \cup \{1, 2\} = \{0, 1, 2\}.
\]
Observamos que $1 \in \partial \Omega \cup \partial \Delta$, pero $1 \notin \partial(\Omega \cup \Delta)$. Por lo tanto, $\partial \Omega \cup \partial \Delta \not\subseteq \partial(\Omega \cup \Delta)$, y la igualdad es falsa.

\textbf{Segunda propiedad:} $\partial(\Omega\cap \Delta) = \partial \Omega\cap \partial \Delta$.

Esta propiedad es \textbf{falsa} en general.

\textit{Contraejemplo:} Tomemos $\Omega = [0,1]$ y $\Delta = [1,2]$ en $\R$. Entonces $\Omega \cap \Delta = \{1\}$. Calculamos las fronteras:
\[
\partial(\Omega \cap \Delta) = \partial(\{1\}) = \{1\},
\]
mientras que
\[
\partial \Omega \cap \partial \Delta = \partial([0,1]) \cap \partial([1,2]) = \{0, 1\} \cap \{1, 2\} = \{1\}.
\]
En este caso, la igualdad sí se cumple. Necesitamos otro contraejemplo.

\textit{Contraejemplo mejor:} Tomemos $\Omega = [0,2]$ y $\Delta = [1,3]$ en $\R$. Entonces $\Omega \cap \Delta = [1,2]$. Calculamos las fronteras:
\[
\partial(\Omega \cap \Delta) = \partial([1,2]) = \{1, 2\},
\]
mientras que
\[
\partial \Omega \cap \partial \Delta = \partial([0,2]) \cap \partial([1,3]) = \{0, 2\} \cap \{1, 3\} = \emptyset.
\]
Claramente, $\partial(\Omega \cap \Delta) = \{1, 2\} \neq \emptyset = \partial \Omega \cap \partial \Delta$. Por lo tanto, la igualdad es falsa.

\item \textbf{Propiedades de monotonía.}

Supongamos que $\Omega \subset \Delta$.

\textbf{Primera propiedad:} $\Int(\Omega)\subset \Int(\Delta)$.

Esta propiedad es \textbf{verdadera}.

\textit{Demostración:} Sea $\x \in \Int(\Omega)$. Entonces existe $\varepsilon > 0$ tal que $B(\x, \varepsilon) \subseteq \Omega$. Como $\Omega \subseteq \Delta$, se tiene que $B(\x, \varepsilon) \subseteq \Delta$. Por lo tanto, $\x \in \Int(\Delta)$. Así, $\Int(\Omega) \subseteq \Int(\Delta)$.

\textbf{Segunda propiedad:} $\Ext(\Omega)\subset \Ext(\Delta)$.

Esta propiedad es \textbf{falsa} en general.

\textit{Contraejemplo:} Tomemos $\Omega = [0,1]$ y $\Delta = [0,2]$ en $\R$. Claramente, $\Omega \subset \Delta$. Calculamos los exteriores:
\[
\Ext(\Omega) = \Ext([0,1]) = \Int(\R \setminus [0,1]) = \Int((-\infty, 0) \cup (1, \infty)) = (-\infty, 0) \cup (1, \infty),
\]
\[
\Ext(\Delta) = \Ext([0,2]) = \Int(\R \setminus [0,2]) = \Int((-\infty, 0) \cup (2, \infty)) = (-\infty, 0) \cup (2, \infty).
\]
Observamos que $1.5 \in \Ext(\Omega)$ (ya que $1.5 \in (1, \infty)$), pero $1.5 \notin \Ext(\Delta)$ (ya que $1.5 \notin (-\infty, 0) \cup (2, \infty)$). Por lo tanto, $\Ext(\Omega) \not\subseteq \Ext(\Delta)$, y la propiedad es falsa.

\textit{Observación:} En realidad, la contención correcta es $\Ext(\Delta) \subseteq \Ext(\Omega)$. Esto se debe a que si $\Omega \subseteq \Delta$, entonces $\R^n \setminus \Delta \subseteq \R^n \setminus \Omega$, y tomando interiores se obtiene $\Ext(\Delta) \subseteq \Ext(\Omega)$.

\textbf{Tercera propiedad:} $\partial \Omega\subset \partial \Delta$.

Esta propiedad es \textbf{falsa} en general.

\textit{Contraejemplo:} Tomemos $\Omega = [0,1]$ y $\Delta = [0,2]$ en $\R$. Claramente, $\Omega \subset \Delta$. Calculamos las fronteras:
\[
\partial \Omega = \partial([0,1]) = \{0, 1\},
\]
\[
\partial \Delta = \partial([0,2]) = \{0, 2\}.
\]
Observamos que $1 \in \partial \Omega$, pero $1 \notin \partial \Delta$. Por lo tanto, $\partial \Omega \not\subseteq \partial \Delta$, y la propiedad es falsa.
\end{enumerate}

\item Define la \emph{suma} de $\Omega$ y $\Delta$ como
\[ \Omega + \Delta := \{ \x+\y\in \R^n : \x\in \Omega, \y\in \Delta \}. \]
Demuestra que $\Omega + \Delta$ es abierto.

\textbf{Solución: } Recordemos que un conjunto $A \subseteq \R^n$ es abierto si para todo punto $\mathbf{p} \in A$, existe $\varepsilon > 0$ tal que la bola abierta $B(\mathbf{p}, \varepsilon)$ está completamente contenida en $A$.

Sea $\mathbf{z} \in \Omega + \Delta$ un punto arbitrario. Por definición de la suma de conjuntos, existen $\x \in \Omega$ y $\y \in \Delta$ tales que $\mathbf{z} = \x + \y$.

Como $\Omega$ es abierto y $\x \in \Omega$, existe $\varepsilon_1 > 0$ tal que
\[
B(\x, \varepsilon_1) \subseteq \Omega.
\]
Análogamente, como $\Delta$ es abierto y $\y \in \Delta$, existe $\varepsilon_2 > 0$ tal que
\[
B(\y, \varepsilon_2) \subseteq \Delta.
\]

Sea $\varepsilon = \min\{\varepsilon_1, \varepsilon_2\} > 0$. Demostraremos que
\[
B(\mathbf{z}, \varepsilon) \subseteq \Omega + \Delta.
\]

Sea $\mathbf{w} \in B(\mathbf{z}, \varepsilon)$ un punto arbitrario. Entonces $\norm{\mathbf{w} - \mathbf{z}} < \varepsilon$. Queremos expresar $\mathbf{w}$ como suma de un elemento de $\Omega$ y un elemento de $\Delta$.

Definimos
\[
\mathbf{u} = \x + (\mathbf{w} - \mathbf{z}) = \x + (\mathbf{w} - \x - \y) = \mathbf{w} - \y.
\]
Observamos que
\[
\norm{\mathbf{u} - \x} = \norm{(\x + \mathbf{w} - \mathbf{z}) - \x} = \norm{\mathbf{w} - \mathbf{z}} < \varepsilon \leq \varepsilon_1.
\]
Por lo tanto, $\mathbf{u} \in B(\x, \varepsilon_1) \subseteq \Omega$.

Ahora, expresamos $\mathbf{w}$ como
\[
\mathbf{w} = \mathbf{u} + \y.
\]
Como $\mathbf{u} \in \Omega$ y $\y \in \Delta$, se sigue que $\mathbf{w} \in \Omega + \Delta$.

Por lo tanto, $B(\mathbf{z}, \varepsilon) \subseteq \Omega + \Delta$.

Como $\mathbf{z} \in \Omega + \Delta$ era arbitrario, hemos demostrado que todo punto de $\Omega + \Delta$ tiene una bola abierta centrada en él completamente contenida en $\Omega + \Delta$. Por lo tanto,
\[
\Omega + \Delta \text{ es abierto.}
\]

\item Encuentra la cerradura de 
\[ \Omega = \{ (x,y)\in \R^2 : x > y^2 \}. \]

\textbf{Solución: } Recordemos que la cerradura de un conjunto $\Omega$, denotada por $\overline{\Omega}$, es el conjunto de todos los puntos de $\Omega$ junto con todos sus puntos de acumulación. Equivalentemente, es el conjunto cerrado más pequeño que contiene a $\Omega$.

Consideremos la función $f(x,y) = x - y^2$. Esta función es continua en todo $\R^2$. El conjunto $\Omega$ está definido por la condición $f(x,y) > 0$, es decir,
\[
\Omega = f^{-1}((0, \infty)).
\]

\textbf{Paso 1: Determinar el interior.}

Dado que $f$ es continua y $(0, \infty)$ es abierto en $\R$, el conjunto $\Omega = f^{-1}((0, \infty))$ es abierto. Por lo tanto,
\[
\Int(\Omega) = \Omega = \{ (x,y)\in \R^2 : x > y^2 \}.
\]

\textbf{Paso 2: Determinar la cerradura.}

La cerradura de $\Omega$ se obtiene ``agregando'' los puntos límite de $\Omega$. Los puntos límite que no están en $\Omega$ son aquellos donde $x = y^2$, es decir, los puntos sobre la parábola $x = y^2$.

Para justificar esto formalmente, demostremos que
\[
\overline{\Omega} = \{ (x,y)\in \R^2 : x \geq y^2 \}.
\]

\textit{Demostración de la contención $\overline{\Omega} \subseteq \{ (x,y)\in \R^2 : x \geq y^2 \}$:}

Sea $A = \{ (x,y)\in \R^2 : x \geq y^2 \}$. Observamos que $A$ es cerrado, ya que $A = f^{-1}([0, \infty))$ y $f$ es continua. Además, $\Omega \subseteq A$, pues si $x > y^2$, entonces ciertamente $x \geq y^2$. Como $\overline{\Omega}$ es el cerrado más pequeño que contiene a $\Omega$, se sigue que
\[
\overline{\Omega} \subseteq A = \{ (x,y)\in \R^2 : x \geq y^2 \}.
\]

\textit{Demostración de la contención $\{ (x,y)\in \R^2 : x \geq y^2 \} \subseteq \overline{\Omega}$:}

Sea $(x_0, y_0) \in A$, es decir, $x_0 \geq y_0^2$. Debemos demostrar que $(x_0, y_0) \in \overline{\Omega}$. Distinguimos dos casos:

\begin{itemize}
    \item \textbf{Caso 1:} $x_0 > y_0^2$. Entonces $(x_0, y_0) \in \Omega$, y por lo tanto $(x_0, y_0) \in \overline{\Omega}$.
    
    \item \textbf{Caso 2:} $x_0 = y_0^2$. Entonces $(x_0, y_0)$ está sobre la parábola. Debemos demostrar que es un punto de acumulación de $\Omega$. Para ello, construiremos una sucesión de puntos en $\Omega$ que converja a $(x_0, y_0)$.
    
    Para cada $k \in \N$, definimos
    \[
    (x_k, y_k) = \left(x_0 + \frac{1}{k}, y_0\right).
    \]
    Observamos que
    \[
    x_k = x_0 + \frac{1}{k} = y_0^2 + \frac{1}{k} > y_0^2 = y_k^2,
    \]
    por lo que $(x_k, y_k) \in \Omega$ para toda $k \in \N$. Además,
    \[
    \lim_{k \to \infty} (x_k, y_k) = \lim_{k \to \infty} \left(x_0 + \frac{1}{k}, y_0\right) = (x_0, y_0).
    \]
    Por lo tanto, $(x_0, y_0)$ es un punto de acumulación de $\Omega$, y así $(x_0, y_0) \in \overline{\Omega}$.
\end{itemize}

En ambos casos, $(x_0, y_0) \in \overline{\Omega}$, lo que demuestra que
\[
\{ (x,y)\in \R^2 : x \geq y^2 \} \subseteq \overline{\Omega}.
\]

Combinando ambas contenciones, concluimos que
\[
\overline{\Omega} = \{ (x,y)\in \R^2 : x \geq y^2 \}.
\]

\item Prueba que si $\Omega$ es cerrado entonces 
$\Int(\partial \Omega)=\emptyset$.

\item Supongamos que $a_1,\ldots,a_n,b_1,\ldots,b_n\in \R^n$ son tales que $a_j<b_j$ para toda $j=1,\ldots,n$. Hemos definimos el rectángulo cerrado (abierto) por
\[ R = [a_1,b_1]\times \cdots\times [a_n,b_n], \]
o bien,
\[ R = (a_1,b_1)\times \cdots\times (a_n,b_n). \]
En cualquier caso, la \textsf{diagonal} de $R$ se define por:
\[ \diag(R) := \norm{(b_1,\ldots,b_n) - (a_1,\ldots,a_n)}. \]
Prueba que todo rectángulo cerrado $R$ es un subconjunto cerrado, acotado y convexo de $\R^n$. 

\end{enumerate}

\end{document}